% !TeX root = main.tex
% PAGINE ORIGINALI 501-510 — DA COMPILARE
% Incollare qui i quiz del blocco e alla fine le rispettive\registraSoluzione{numero}{lettera}.
% Questo file vuoto non aggiunge pagine al PDF.

% !TeX root = main.tex
% Pagine 501-510: spiegazioni 30a simulazione, test di recupero 1 e primi 20 quiz del test 2.

% Pag. originale 502/521
\iffalse Pag. originale 502/521 \fi
\begin{quiz}
Il dominio della funzione $f(x)=\log\frac{(x^2+1)(4-x^2)}{x^2-2x+1}$ \`e:
\begin{enumerate}
\item $(-2,2)$
\item $(-2,1)\cup(1,2)$
\item $(-\infty,-2)\cup(2,+\infty)$
\item $[-2,1)\cup(1,2]$
\item $\mathbb R\setminus\{1\}$
\end{enumerate}
\end{quiz}
% SOLUZIONE: B
% SEZIONE_AM1: 1.4
% SPIEGAZIONE: Il numeratore ha il segno di $4-x^2$, mentre il denominatore $(x-1)^2$ \`e positivo per $x\ne1$. Il logaritmo richiede $-2<x<2$, con $x\ne1$.

% Pag. originale 502/521
\iffalse Pag. originale 502/521 \fi
\begin{quiz}
La funzione $f:\mathbb R\to\mathbb R$, $f(x)=\sqrt{|x+1|-2x}$
\begin{enumerate}
\item \`e suriettiva
\item non \`e invertibile
\item ha dominio $(-\infty,1)$
\item $\operatorname{im}f=(-1,+\infty)$
\item \`e iniettiva
\end{enumerate}
\end{quiz}
% SOLUZIONE: E
% SEZIONE_AM1: 1.4
% SPIEGAZIONE: Il radicando \`e $-3x-1$ per $x<-1$ e $1-x$ per $x\ge-1$. Nel dominio $(-\infty,1]$ la funzione \`e strettamente decrescente, quindi iniettiva. Il testo scrive $f:\mathbb R\to\mathbb R$, ma la formula non \`e definita per $x>1$: incongruenza nell\'indicazione del dominio.

% Pag. originale 502/521
\iffalse Pag. originale 502/521 \fi
\begin{quiz}
\`E data l\'equazione di variabile complessa $\left|\overline z-\frac9z\right|(z^3-8i)=0$. Quale delle seguenti affermazioni \`e corretta?
\begin{enumerate}
\item L\'equazione ammette soltanto $3$ soluzioni complesse
\item $z=-2i$ \`e la sola soluzione immaginaria dell\'equazione
\item Tra le infinite soluzioni dell\'equazione solo due sono reali
\item Le soluzioni dell\'equazione, rappresentate sul piano di Gauss, sono i punti che distano dall\'origine $3$
\item Sono soluzioni dell\'equazione i due numeri complessi $z_{1,2}=\sqrt3\pm i$
\end{enumerate}
\end{quiz}
% SOLUZIONE: C
% SEZIONE_AM1: 1.5
% SPIEGAZIONE: La condizione $\overline z=9/z$ equivale a $|z|=3$, che fornisce una circonferenza di soluzioni. Anche $z^3=8i$ fornisce tre radici di modulo $2$. Le uniche soluzioni reali sono $z=\pm3$.

% Pag. originale 502/521
\iffalse Pag. originale 502/521 \fi
\begin{quiz}
$\displaystyle\lim_{x\to0^+}\frac{\sqrt{x^4+x^6}-x^2}{x^4}$ vale
\begin{enumerate}
\item $+\infty$
\item $0$
\item $1$
\item $2$
\item $\frac12$
\end{enumerate}
\end{quiz}
% SOLUZIONE: E
% SEZIONE_AM1: 3.4
% SPIEGAZIONE: Per $x>0$, $\sqrt{x^4+x^6}=x^2\sqrt{1+x^2}$; razionalizzando si ottiene $1/(\sqrt{1+x^2}+1)\to1/2$.

% Pag. originale 502/521
\iffalse Pag. originale 502/521 \fi
\begin{quiz}
Se $f=o(g)$ e $g=o(x^5)$, per $x\to0$, allora, per $x\to0$, si pu\`o concludere che
\begin{enumerate}
\item $f(x)g(x)=o(x^{11})$
\item $f(x)g(x)\sim x^{10}$
\item $f(x)+x^5=o(g(x))$
\item $\frac{f(x)}{g(x)}\to0$
\item nessuna delle risposte precedenti
\end{enumerate}
\end{quiz}
% SOLUZIONE: D
% SEZIONE_AM1: 3.1
% SPIEGAZIONE: $f=o(g)$ significa esattamente $f/g\to0$ dove il rapporto \`e definito. Il prodotto \`e soltanto $o(x^{10})$, non necessariamente $o(x^{11})$.

% Pag. originale 502/521
\iffalse Pag. originale 502/521 \fi
\begin{quiz}
La funzione $f(x)=\log\left(3+\frac4{x^2}\right)+2x$ ha come asintoto obliquo sinistro la retta:
\begin{enumerate}
\item $y=3x+\log3$
\item $y=2x+\log4$
\item $y=5x+\log4$
\item $y=2x$
\item $y=2x+\log3$
\end{enumerate}
\end{quiz}
% SOLUZIONE: E
% SEZIONE_AM1: 4.4
% SPIEGAZIONE: Per $x\to-\infty$, $\log(3+4/x^2)\to\log3$; dunque $f(x)-(2x+\log3)\to0$.

% Pag. originale 503/521
\iffalse Pag. originale 503/521 \fi
\begin{quiz}
$\displaystyle\lim_{n\to+\infty}\frac{\log n^{10}+\sin n}{\sqrt n-3}=$
\begin{enumerate}
\item $10$
\item non esiste
\item $+\infty$
\item $0$
\item $1$
\end{enumerate}
\end{quiz}
% SOLUZIONE: D
% SEZIONE_AM1: 3.1
% SPIEGAZIONE: Il numeratore \`e $10\log n+O(1)$, il denominatore \`e equivalente a $\sqrt n$; il rapporto tende a zero.

% Pag. originale 503/521
\iffalse Pag. originale 503/521 \fi
\begin{quiz}
Sia data la funzione $f(x)=\log|\sin2x|$. La sua derivata prima vale:
\begin{enumerate}
\item $f\prime(x)=\frac{\cos2x}{\sin2x}$
\item $f\prime(x)=2\frac{\cos x}{\sin x}$
\item $f\prime(x)=\frac1{|\sin2x|}$
\item $f\prime(x)=2\frac{\cos2x}{\sin2x}$
\item $f\prime(x)=\frac2{\sin2x}$
\end{enumerate}
\end{quiz}
% SOLUZIONE: D
% SEZIONE_AM1: 4.1
% SPIEGAZIONE: La derivata di $\log|u|$ \`e $u\prime/u$: con $u=\sin2x$ si ottiene $2\cos2x/\sin2x$.

% Pag. originale 503/521
\iffalse Pag. originale 503/521 \fi
\begin{quiz}
Sia $f:\mathbb R\to\mathbb R$ una funzione derivabile tale che $f(0)=1$, $f(4)=e$, $f\prime(0)=2$. Si consideri la funzione $h(x)=\log f(x)$; quale delle seguenti risposte \`e l\'unica esatta?
\begin{enumerate}
\item Esiste almeno un punto $c\in(0,4)$ tale che $h\prime(c)=\frac{e-1}4$
\item $h\prime(0)=2$
\item Esiste almeno un punto $c\in(0,4)$ tale che $h\prime(c)=-\frac14$
\item Esiste almeno un punto $c\in(0,4)$ tale che $h(c)=2$
\item Esiste almeno un punto $c\in(0,4)$ tale che $h\prime(c)=\frac4{e-1}$
\end{enumerate}
\end{quiz}
% SOLUZIONE: B
% SEZIONE_AM1: 4.2
% SPIEGAZIONE: Poich\'e $f(0)=1>0$, $\log f$ \`e definita vicino a zero e $h\prime(0)=f\prime(0)/f(0)=2$. Non \`e garantita la positivit\`a di $f$ su tutto $[0,4]$.

% Pag. originale 503/521
\iffalse Pag. originale 503/521 \fi
\begin{quiz}
Sia $f:\mathbb R\to\mathbb R$ una funzione derivabile tale che $f(0)=0$, $f(3)=2$. Considerando $g(x)=e^{-f(x)}$, allora esiste un punto $c\in(0,3)$ tale che
\begin{enumerate}
\item $g\prime(c)=\frac{1-e^{-2}}3$
\item $g\prime(c)=\frac23$
\item $g\prime(c)=\frac{1-e^2}3$
\item $g\prime(c)=\frac{e^2-1}{3e^2}$
\item $g\prime(c)=\frac{e^{-2}-1}3$
\end{enumerate}
\end{quiz}
% SOLUZIONE: E
% SEZIONE_AM1: 4.2
% SPIEGAZIONE: Per Lagrange $g\prime(c)=[g(3)-g(0)]/3=(e^{-2}-1)/3$.

% Pag. originale 503/521
\iffalse Pag. originale 503/521 \fi
\begin{quiz}
Lo sviluppo di Mac Laurin, di ordine $4$, della funzione $f(x)=\sin^2(x-x^2)-x^2+2x^3$ \`e
\begin{enumerate}
\item $\frac12x^4+o(x^4)$
\item $\frac23x^4$
\item $\frac23x^4+o(x^4)$
\item $4x^3+\frac23x^4+o(x^4)$
\item nessuno dei precedenti
\end{enumerate}
\end{quiz}
% SOLUZIONE: C
% SEZIONE_AM1: 3.4
% SPIEGAZIONE: Usando $\sin^2t=t^2-t^4/3+o(t^4)$ con $t=x-x^2$, si trova $x^2-2x^3+(1-1/3)x^4+o(x^4)$; i termini fino al cubico si cancellano.

% Pag. originale 504/521
\iffalse Pag. originale 504/521 \fi
\begin{quiz}
Sia $f$ una funzione di classe $C^{(3)}(\mathbb R)$, tale che $f(x)=-2+2x^2-3x^3+o(x^3)$. \`E vero che:
\begin{enumerate}
\item $f(1)=-3$
\item $f(0)=-2$ e $f\prime(0)=0$
\item $f(0)=-2$ e $f\prime\prime(0)=2$
\item Il punto $x=0$ \`e un punto di massimo relativo
\item $f\prime\prime\prime(0)=-3$
\end{enumerate}
\end{quiz}
% SOLUZIONE: B
% SEZIONE_AM1: 3.4
% SPIEGAZIONE: Confrontando con Taylor si ha $f(0)=-2$, $f\prime(0)=0$, $f\prime\prime(0)=4$, $f\prime\prime\prime(0)=-18$. Inoltre $x=0$ \`e un minimo locale.

% Pag. originale 504/521
\iffalse Pag. originale 504/521 \fi
\begin{quiz}
L\'enunciato corretto \`e:
\begin{enumerate}
\item se $f$ \`e definita su $[a,b]$ e $f(a)f(b)<0$, allora $f$ ha almeno uno zero in $(a,b)$
\item se $f$ \`e continua su $[a,b]$ e $f(a)f(b)<0$, allora $f$ ha almeno uno zero in $(a,b)$
\item se $f$ \`e continua su $(a,b)$ e $f(a)f(b)<0$, allora $f$ ha almeno uno zero in $(a,b)$
\item se $f$ \`e continua su $[a,b]$ e $f(a)f(b)<0$, allora $f$ ha al pi\`u uno zero in $(a,b)$
\item se $f$ \`e continua su $[a,b]$ e $f(a)f(b)<0$, allora $0\in(a,b)$
\end{enumerate}
\end{quiz}
% SOLUZIONE: B
% SEZIONE_AM1: 3.2
% SPIEGAZIONE: Il teorema degli zeri richiede continuit\`a sull\'intervallo chiuso e valori agli estremi di segno opposto.

% Pag. originale 504/521
\iffalse Pag. originale 504/521 \fi
\begin{quiz}
Sia $f$ una funzione continua in un intervallo $I$. Una primitiva di $f$ su $I$ \`e una funzione $F$ tale che:
\begin{enumerate}
\item $F(x)=\int_{x_0}^x f(t)\,dt$, con $x_0,x\in I$
\item $F$ \`e definita su $I$ e $F\prime(x)=f(x)$, $\forall x\in I$
\item $F$ \`e derivabile su $I$ e $F\prime(x)=f(x)$, $\forall x\in I$
\item $f\prime(x)=F(x)$, $\forall x\in I$
\item $F$ \`e derivabile in $x_0\in I$ e $F\prime(x_0)=f(x_0)$
\end{enumerate}
\end{quiz}
% SOLUZIONE: C
% SEZIONE_AM1: 5.2
% SPIEGAZIONE: La definizione di primitiva richiede derivabilit\`a su tutto $I$ e $F\prime=f$ su $I$. La (a) ne d\`a una costruzione particolare, non tutte le primitive.

% Pag. originale 504/521
\iffalse Pag. originale 504/521 \fi
\begin{quiz}
Sia data la funzione $f(x)=4\log(\cosh x)-2x^2-3x^3$. La parte principale della funzione $f$, per $x\to0$, rispetto all\'infinitesimo campione standard $u(x)$, \`e:
\begin{enumerate}
\item $-3x^3$
\item $-2x^2$
\item $-3x^3+o(x^3)$
\item $-2x^2-3x^3$
\item $-3x^3-\frac{x^4}3$
\end{enumerate}
\end{quiz}
% SOLUZIONE: A
% SEZIONE_AM1: 3.4
% SPIEGAZIONE: $\log\cosh x=x^2/2-x^4/12+O(x^6)$, quindi $4\log\cosh x-2x^2-3x^3=-3x^3-x^4/3+O(x^6)$; parte principale $-3x^3$.

% Pag. originale 504/521
\iffalse Pag. originale 504/521 \fi
\begin{quiz}
Il seguente integrale definito $\displaystyle\int_2^4\frac{dx}{x(2-\log x)}$ vale:
\begin{enumerate}
\item $\log\frac{2-\log2}{\ln2-1}-\log2$
\item $\log\frac{2-\log2}{1-\ln2}$
\item $\log\frac{2-\log2}{2\ln2-2}$
\item $|\log(2-\log2)-\log2|$
\item $\log\frac{2-\log2}{1-\ln2}-\log2$
\end{enumerate}
\end{quiz}
% SOLUZIONE: E
% SEZIONE_AM1: 5.4
% SPIEGAZIONE: Con $u=2-\log x$, l\'integrale \`e $\log[(2-\log2)/(2-\log4)]$. Poich\'e $2-\log4=2(1-\log2)$, si ottiene la (e).

% Pag. originale 505/521
\iffalse Pag. originale 505/521 \fi
\begin{quiz}
Qual \`e la corretta formulazione del teorema della media integrale?
\begin{enumerate}
\item Sia $f$ integrabile in $[a,b]$. La media integrale $\mu$ soddisfa $\min_{x\in[a,b]}f(x)\le\mu=f(c)\le\max_{x\in[a,b]}f(x)$, $c\in[a,b]$
\item Sia $f$ integrabile in $[a,b]$. La media integrale $\mu$ soddisfa $\min_{x\in[a,b]}f(x)\le\mu\le\max_{x\in[a,b]}f(x)$
\item Sia $f$ integrabile in $[a,b]$. La media integrale $\mu$ soddisfa $\inf_{x\in[a,b]}f(x)\le\mu\le\sup_{x\in[a,b]}f(x)$
\item Sia $f$ continua in $[a,b]$ ed integrabile almeno in $(a,b)$. Detta $\mu$ la media integrale, $\mu=\frac1{b-a}\int_a^b f\prime(x)\,dx$
\item Sia $f$ derivabile in $[a,b]$. La media integrale $\mu$ soddisfa $\min_{x\in[a,b]}f(x)\le\mu=f\prime(c)\le\max_{x\in[a,b]}f(x)$, $c\in[a,b]$
\end{enumerate}
\end{quiz}
% SOLUZIONE: C
% SEZIONE_AM1: 5.2
% SPIEGAZIONE: Per una funzione integrabile il valor medio \`e compreso tra infimo e supremo, anche se il minimo o massimo non esistono. L\'uguaglianza $\mu=f(c)$ richiede ulteriori ipotesi, ad esempio continuit\`a.

% Pag. originale 505/521
\iffalse Pag. originale 505/521 \fi
\begin{quiz}
Sia $f:\mathbb R\to\mathbb R$ continua e sia $F(x)=\int_x^0 f(t)\,dt$. Se $f\ge0$ per ogni $x\in\mathbb R$, allora si pu\`o dedurre che:
\begin{enumerate}
\item $F(x)\le0$, $\forall x\in\mathbb R$
\item $F(x)\ge0$, $\forall x\in\mathbb R$
\item $F(x)$ \`e decrescente in $\mathbb R$
\item $F(x)\ge0$, $\forall x\in(0,+\infty)$
\item $F(x)$ \`e crescente $\forall x\in(-\infty,0)$
\end{enumerate}
\end{quiz}
% SOLUZIONE: C
% SEZIONE_AM1: 5.2
% SPIEGAZIONE: Per il teorema fondamentale del calcolo $F\prime(x)=-f(x)\le0$, quindi $F$ \`e non crescente.

% Pag. originale 505/521
\iffalse Pag. originale 505/521 \fi
\begin{quiz}
Quale delle seguenti propriet\`a \`e soddisfatta dalla funzione $f(x)=\frac{\sin x}{x\sqrt x}$?
\begin{enumerate}
\item $\int_0^{+\infty}f(x)\,dx$ converge
\item $\int_0^\pi f(x)\,dx$ diverge
\item $\int_1^{+\infty}f(x)\,dx$ non converge assolutamente
\item $\int_\pi^{+\infty}|f(x)|\,dx$ diverge
\item $f(x)\sim1/\sqrt x$ per $x\to0$; $f(x)\sim1/\sqrt x$ per $x\to+\infty$
\end{enumerate}
\end{quiz}
% SOLUZIONE: A
% SEZIONE_AM1: 5.5
% SPIEGAZIONE: Vicino a zero $\sin x/(x\sqrt x)\sim1/\sqrt x$, integrabile. All\'infinito $|f(x)|\le x^{-3/2}$, anch\'esso integrabile: converge assolutamente su $(0,+\infty)$.

% Pag. originale 505/521
\iffalse Pag. originale 505/521 \fi
\begin{quiz}
L\'equazione differenziale $x\prime(t)=(2-x)(1-\sin x)(\pi-t)\sqrt x$
\begin{enumerate}
\item non \`e a variabili separabili
\item \`e lineare
\item ammette la soluzione costante $t=\pi$
\item ammette le soluzioni costanti $x(t)=2$, $x(t)=0$, $x(t)=\frac\pi2+2k\pi$, $k\in\mathbb N$, $\forall t\in\mathbb R$
\item ammette le soluzioni costanti $t=\pi$, $x(t)=2$, $x(t)=0$, $x(t)=\frac\pi2+2k\pi$, $k\in\mathbb Z$, $\forall t\in\mathbb R$
\end{enumerate}
\end{quiz}
% SOLUZIONE: D
% SEZIONE_AM1: 6.1
% SPIEGAZIONE: Annullando il fattore dipendente da $x$ si ottengono le soluzioni costanti $x=0$, $x=2$ e $x=\pi/2+2k\pi$ con $x\ge0$. La risposta (d) segue la convenzione $\mathbb N=\{0,1,\dots\}$.

% Pag. originale 507/521
\iffalse Pag. originale 507/521 \fi
\begin{quiz}
Il dominio della funzione $f(x)=\log(\sqrt{x-2}-3)$ \`e:
\begin{enumerate}
\item $[2,3]$
\item $[11,+\infty)$
\item $(5,+\infty)$
\item $(11,+\infty)$
\item $[3,+\infty)$
\end{enumerate}
\end{quiz}
% SOLUZIONE: D
% SEZIONE_AM1: 1.4
% SPIEGAZIONE: Per definire il logaritmo occorre $\sqrt{x-2}>3$, cio\`e $x>11$.

% Pag. originale 507/521
\iffalse Pag. originale 507/521 \fi
\begin{quiz}
La funzione $f(x)=|\log(x-3)|$
\begin{enumerate}
\item \`e invertibile sull\'intervallo $[3,+\infty)$
\item \`e invertibile sull\'intervallo $[4,+\infty)$
\item \`e invertibile sull\'intervallo $(0,+\infty)$
\item non \`e invertibile su alcun intervallo
\item \`e invertibile sull\'intervallo $(3,5)$
\end{enumerate}
\end{quiz}
% SOLUZIONE: B
% SEZIONE_AM1: 1.4
% SPIEGAZIONE: Su $[4,+\infty)$ si ha $\log(x-3)\ge0$ e la funzione \`e crescente, quindi iniettiva. Negli intervalli che attraversano $x=4$ il valore assoluto rompe la monotonia.

% Pag. originale 507/521
\iffalse Pag. originale 507/521 \fi
\begin{quiz}
Sia $z=1-2i$. Allora $|z^2+\overline z|$ vale
\begin{enumerate}
\item $2\sqrt2$
\item $8$
\item $4$
\item $4\sqrt2$
\item $2$
\end{enumerate}
\end{quiz}
% SOLUZIONE: A
% SEZIONE_AM1: 1.5
% SPIEGAZIONE: $z^2=-3-4i$ e $\overline z=1+2i$; somma $-2-2i$, modulo $2\sqrt2$.

% Pag. originale 507/521
\iffalse Pag. originale 507/521 \fi
\begin{quiz}
Il limite $\displaystyle\lim_{x\to+\infty}\frac{e^{-2x}-2x+\cos x}{e^{-x}+3x-3\sin x}$ vale
\begin{enumerate}
\item $-2$
\item $-1$
\item $2$
\item $0$
\item $-\frac23$
\end{enumerate}
\end{quiz}
% SOLUZIONE: E
% SEZIONE_AM1: 3.1
% SPIEGAZIONE: Gli esponenziali decadono e i termini trigonometrici restano limitati: il rapporto \`e equivalente a $-2x/(3x)=-2/3$.

% Pag. originale 507/521
\iffalse Pag. originale 507/521 \fi
\begin{quiz}
Per $x\to0$, sia $f(x)\sim x^2\cos x$ e $g(x)\sim e^x-1$. Allora si ha
\begin{enumerate}
\item $\lim_{x\to0}f(x)/g(x)=+\infty$
\item $\lim_{x\to0}f(x)/g(x)=\frac12$
\item $\lim_{x\to0}f(x)/g(x)=0$
\item $\lim_{x\to0}f(x)/g(x)=+\infty$
\item $\lim_{x\to0}f(x)/g(x)=1$
\end{enumerate}
\end{quiz}
% SOLUZIONE: C
% SEZIONE_AM1: 3.1
% SPIEGAZIONE: $f(x)\sim x^2$ e $g(x)\sim x$, dunque $f/g\sim x\to0$. Le opzioni (a) e (d) sono duplicate nel testo originale.

% Pag. originale 508/521
\iffalse Pag. originale 508/521 \fi
\begin{quiz}
Il limite $\displaystyle\lim_{x\to+\infty}(3x^3-4x^2)\sin x$ vale
\begin{enumerate}
\item $-\infty$
\item $+\infty$
\item $0$
\item sia $+\infty$ che $-\infty$
\item non esiste
\end{enumerate}
\end{quiz}
% SOLUZIONE: E
% SEZIONE_AM1: 3.1
% SPIEGAZIONE: L\'ampiezza cresce senza limite e $\sin x$ oscilla: lungo sottosuccessioni il prodotto diverge a $+\infty$ e a $-\infty$. Pertanto il limite non esiste; la (d) descrive impropriamente due limiti diversi.

% Pag. originale 508/521
\iffalse Pag. originale 508/521 \fi
\begin{quiz}
Sia $a_n$ una successione limitata inferiormente. Allora
\begin{enumerate}
\item $\forall k>0\ \exists\overline n\in\mathbb N$ tale che $n>\overline n\Rightarrow a_n\ge k$
\item per ogni $n$, $a_n\ge0$
\item $\exists k<0$ tale che per ogni $n$, $a_n<k$
\item $\exists\overline n\in\mathbb N$ tale che $n>\overline n\Rightarrow a_n\ge0$
\item $\exists k<0$ tale che per ogni $n$, $a_n>k$
\end{enumerate}
\end{quiz}
% SOLUZIONE: E
% SEZIONE_AM1: 2.1
% SPIEGAZIONE: Una successione limitata inferiormente ammette un minorante reale $m$; scegliendo $k<\min\{m,0\}$ si ha $a_n>k$ per ogni $n$.

% Pag. originale 508/521
\iffalse Pag. originale 508/521 \fi
\begin{quiz}
Sia $f(x)=(3\cos x)^{4\cos x}$. Allora $f\prime(0)$ vale
\begin{enumerate}
\item $1$
\item $0$
\item $\frac34$
\item $\log\frac34$
\item $-1$
\end{enumerate}
\end{quiz}
% SOLUZIONE: B
% SEZIONE_AM1: 4.1
% SPIEGAZIONE: La derivazione logaritmica d\`a $f\prime/f=-4\sin x\log(3\cos x)-4\sin x$; in $x=0$ entrambi i termini sono nulli.

% Pag. originale 508/521
\iffalse Pag. originale 508/521 \fi
\begin{quiz}
La funzione $f(x)=\begin{cases}2x&\text{se }x\le0,\\2\sqrt x&\text{se }x>0\end{cases}$
\begin{enumerate}
\item \`e di classe $C^1$
\item \`e derivabile nell\'origine
\item \`e discontinua nell\'origine
\item \`e continua ma non \`e derivabile nell\'origine
\item nessuna delle altre affermazioni \`e corretta
\end{enumerate}
\end{quiz}
% SOLUZIONE: D
% SEZIONE_AM1: 4.1
% SPIEGAZIONE: Entrambi i rami tendono a zero; la derivata destra del ramo $2\sqrt x$ tende a $+\infty$, mentre quella sinistra vale $2$.

% Pag. originale 508/521
\iffalse Pag. originale 508/521 \fi
\begin{quiz}
Sia $f:\mathbb R\to\mathbb R$ una funzione derivabile tale che $f(0)=4$, $f\prime(0)=3$. Considerando $h(x)=\frac1{f(x)}$, risulta
\begin{enumerate}
\item $h\prime(0)=-\frac3{16}$
\item $h\prime(0)=-\frac43$
\item $h\prime(0)=\frac13$
\item $h\prime(0)=\frac3{16}$
\item $h\prime(0)=-\frac13$
\end{enumerate}
\end{quiz}
% SOLUZIONE: A
% SEZIONE_AM1: 4.1
% SPIEGAZIONE: $h\prime=-f\prime/f^2$, quindi $h\prime(0)=-3/16$.

% Pag. originale 509/521
\iffalse Pag. originale 509/521 \fi
\begin{quiz}
Il polinomio di McLaurin di ordine $6$ della funzione $f(x)=e^{\cos(x^3)}$ \`e
\begin{enumerate}
\item $1+\frac12x^6$
\item $e-\frac e2x^6$
\item $2+\frac12x^6$
\item $1+\frac16x^6$
\item $1+\frac16x^6$
\end{enumerate}
\end{quiz}
% SOLUZIONE: B
% SEZIONE_AM1: 3.4
% SPIEGAZIONE: $\cos(x^3)=1-x^6/2+O(x^{12})$; quindi $e^{\cos(x^3)}=e(1-x^6/2+O(x^{12}))$. Le alternative (d) ed (e) risultano identiche nel PDF.

% Pag. originale 509/521
\iffalse Pag. originale 509/521 \fi
\begin{quiz}
Se $f$ ha sviluppo di Taylor $f(x)=4-3(x-2)^6+o((x-2)^6)$ per $x\to2$, allora
\begin{enumerate}
\item $f$ ha un flesso in $x=2$
\item $f$ ha un massimo in $x=0$
\item $f$ ha un minimo in $x=0$
\item $f$ ha un minimo in $x=2$
\item $f$ ha un massimo in $x=2$
\end{enumerate}
\end{quiz}
% SOLUZIONE: E
% SEZIONE_AM1: 3.4
% SPIEGAZIONE: Per $x$ vicino a $2$ e diverso da $2$, $f(x)-f(2)=-3(x-2)^6+o((x-2)^6)<0$: massimo locale stretto.

% Pag. originale 509/521
\iffalse Pag. originale 509/521 \fi
\begin{quiz}
Sia $f:[0,3]\to\mathbb R$ una funzione continua e decrescente. Allora si pu\`o dedurre che
\begin{enumerate}
\item $f([0,3])=[f(0),f(3)]$
\item $f((0,3))=\{f(3),f(0)\}$
\item $f((0,3])=(f(3),f(0)]$
\item $f([0,3])=[f(3),f(0)]$
\item nessuna delle altre risposte \`e corretta
\end{enumerate}
\end{quiz}
% SOLUZIONE: D
% SEZIONE_AM1: 1.4
% SPIEGAZIONE: La continuit\`a garantisce tutti i valori intermedi e la decrescenza individua il massimo in $0$ e il minimo in $3$.

% Pag. originale 509/521
\iffalse Pag. originale 509/521 \fi
\begin{quiz}
Sia $f:\mathbb R\to\mathbb R$ una funzione derivabile tale che $f(0)=f(1)=0$. Ponendo $g(x)=f^4(x)$, allora
\begin{enumerate}
\item la derivata $g\prime(x)$ si annulla almeno tre volte
\item la derivata $g\prime(x)$ si annulla esattamente due volte
\item la derivata $g\prime(x)$ si annulla esattamente tre volte
\item la derivata $g\prime(x)$ non si annulla mai
\item la derivata $g\prime(x)$ si annulla almeno quattro volte
\end{enumerate}
\end{quiz}
% SOLUZIONE: A
% SEZIONE_AM1: 4.2
% SPIEGAZIONE: $g\prime(x)=4f(x)^3f\prime(x)$, quindi si annulla in $0$ e in $1$. Per Rolle applicato a $f$ esiste inoltre $c\in(0,1)$ con $f\prime(c)=0$: almeno tre zeri distinti.

% Pag. originale 509/521
\iffalse Pag. originale 509/521 \fi
\begin{quiz}
La parte principale (rispetto a $\varphi(x)=x$) per $x\to0^+$ di $f(x)=3x+\sqrt{4x^2+2x^3}$ \`e
\begin{enumerate}
\item $5x$
\item $x^{3/2}$
\item $\sqrt2x^{3/2}$
\item $3x$
\item $3x+o(x)$
\end{enumerate}
\end{quiz}
% SOLUZIONE: A
% SEZIONE_AM1: 3.4
% SPIEGAZIONE: Per $x>0$, $\sqrt{4x^2+2x^3}=x\sqrt{4+2x}\sim2x$; quindi $f(x)\sim5x$.

% Pag. originale 509/521
\iffalse Pag. originale 509/521 \fi
\begin{quiz}
Una primitiva della funzione $f(x)=\frac{3x}{2x^2+2}$ \`e:
\begin{enumerate}
\item $\frac34\log(x^2+1)$
\item $\frac34\arctan x$
\item $3\log(x^2+1)$
\item $\frac34\arctan(x^2+1)$
\item nessuna delle altre risposte
\end{enumerate}
\end{quiz}
% SOLUZIONE: A
% SEZIONE_AM1: 5.4
% SPIEGAZIONE: $f(x)=(3/2)x/(x^2+1)$; integrando con $u=x^2+1$ si ottiene $(3/4)\log(x^2+1)+C$.

% Pag. originale 510/521
\iffalse Pag. originale 510/521 \fi
\begin{quiz}
Dire quale tra i seguenti enunciati \`e corretto.
\begin{enumerate}
\item se $f$ \`e integrabile in $[a,b]$, allora $\exists c\in[a,b]$ tale che $f(c)=\frac1{b-a}\int_a^b f(x)\,dx$
\item se $f$ \`e continua in $[a,b]$, allora $\exists c\in[a,b]$ tale che $f\prime(c)=\frac1{b-a}\int_a^b f(x)\,dx$
\item se $f$ \`e continua in $[a,b]$, allora $\exists c\in[a,b]$ tale che $f(c)=\frac1{b-a}\int_a^b f(x)\,dx$
\item se $f$ \`e continua in $[a,b]$, allora $\exists c\in[a,b]$ tale che $f(c)=\int_a^b f(x)\,dx$
\item se $f$ \`e integrabile in $[a,b]$, allora $\exists c\in[a,b]$ tale che $f(c)=\int_a^b f(x)\,dx$
\end{enumerate}
\end{quiz}
% SOLUZIONE: C
% SEZIONE_AM1: 5.2
% SPIEGAZIONE: Il teorema della media integrale nella forma con $f(c)=\mu$ richiede la continuit\`a su $[a,b]$.

% Pag. originale 510/521
\iffalse Pag. originale 510/521 \fi
\begin{quiz}
Sia $F(x)=\int_0^x t^2\cosh(t^2)\,dt$. Allora
\begin{enumerate}
\item $F$ \`e crescente su $(0,+\infty)$ e decrescente su $(-\infty,0)$
\item $F$ \`e crescente su $\mathbb R$
\item $F$ ha un minimo in $0$
\item $F$ ha un massimo in $0$
\item nessuna delle altre risposte
\end{enumerate}
\end{quiz}
% SOLUZIONE: B
% SEZIONE_AM1: 5.2
% SPIEGAZIONE: $F\prime(x)=x^2\cosh(x^2)\ge0$ ed \`e positiva per $x\ne0$: $F$ \`e strettamente crescente su $\mathbb R$.

% Pag. originale 510/521
\iffalse Pag. originale 510/521 \fi
\begin{quiz}
Sia $f$ continua su $[0,+\infty)$ e tale che $f(x)\le0$ per ogni $x\ge0$. Allora necessariamente l\'integrale improprio $\int_0^{+\infty}f(x)\,dx$
\begin{enumerate}
\item \`e indeterminato
\item diverge a $-\infty$
\item converge a un numero negativo
\item converge, oppure diverge a $-\infty$
\item nessuna delle precedenti risposte
\end{enumerate}
\end{quiz}
% SOLUZIONE: D
% SEZIONE_AM1: 5.5
% SPIEGAZIONE: Le primitive parziali $\int_0^R f$ sono non crescenti in $R$; il loro limite esteso esiste ed \`e finito non positivo oppure $-\infty$.

% Pag. originale 510/521
\iffalse Pag. originale 510/521 \fi
\begin{quiz}
L\'equazione differenziale $y\prime\prime-y=0$
\begin{enumerate}
\item ha almeno una soluzione non limitata su $(0,+\infty)$
\item non ha soluzioni limitate su $(0,+\infty)$
\item non ha soluzioni illimitate su $(0,+\infty)$
\item ha almeno una soluzione che cambia segno infinite volte
\item ha solo soluzioni positive
\end{enumerate}
\end{quiz}
% SOLUZIONE: A
% SEZIONE_AM1: 6.1
% SPIEGAZIONE: Le soluzioni sono $y=C_1e^x+C_2e^{-x}$: scegliendo $C_1=1$, $C_2=0$ si ha una soluzione illimitata; scegliendo $C_1=0$ si hanno soluzioni limitate.
Risposta completata
